\section{Architecture \& Implementation}

The \texttt{FloorSinhActivation} module is implemented as a lightweight PyTorch \texttt{nn.Module}.

\subsection{PyTorch Production Code}

\begin{lstlisting}[language=Python]
import torch
import torch.nn as nn

class FloorSinhActivation(nn.Module):
    """Enhancement #28: Floor-Sinh Activation Regularizer."""

    def __init__(self, physical_floor: float = -1.0):
        super().__init__()
        self.floor = physical_floor

    def forward(self, x: torch.Tensor) -> torch.Tensor:
        sinh_tensor = torch.sinh(x)
        regularized_tensor = torch.clamp(sinh_tensor, min=self.floor)
        return regularized_tensor
\end{lstlisting}

\subsection{Computational Graph Optimization}
\begin{enumerate}
    \item \textbf{Hyperbolic Pass}: `torch.sinh(x)` is evaluated in a single vector kernel on GPU/CPU.
    \item \textbf{Tensor Clamping}: `torch.clamp(..., min=self.floor)` applies elementwise lower bounding in place or via fused CUDA ops.
    \item \textbf{Feedforward Integration}: Functions as a drop-in replacement for `nn.ReLU()` or `nn.GELU()` in NRC transformer feedforward blocks.
\end{enumerate}
